\documentclass[11pt]{article}
\usepackage[margin=1in]{geometry}
\usepackage{amsmath,amssymb}
\usepackage[hidelinks]{hyperref}
\setlength{\parindent}{0pt}
\setlength{\parskip}{0.7em}

\title{Literature Status: Nystr\"om-SGM Beyond the Saturation Threshold}
\author{Run \texttt{fa\_banach\_001}, lane 2}
\date{11 August 2026}

\begin{document}
\maketitle

\textbf{Status.} \texttt{literature\_already\_answered}.  This is an exact
later-literature match, not a new result of the run.

\section*{Source question}

Junhong Lin and Lorenzo Rosasco, \emph{Optimal Rates for Learning with
Nystr\"om Stochastic Gradient Methods}, arXiv:1710.07797v1 (2017), Section~4,
PDF p.~10, state that their NySGM bounds saturate at the source exponent
$r=1/2$ and conjecture:

\begin{quote}
``by considering a different subsampling technique, we may get optimal error
bounds even for $r\geq 1/2$.''
\end{quote}

Their optimal prediction-error rate, in the notation used there, is
\[
 n^{-\frac{2r+1}{2r+\gamma+1}}
\]
up to logarithmic factors.

\section*{Supporting answer}

Junhong Lin and Volkan Cevher, \emph{Convergences of Regularized Algorithms
and Stochastic Gradient Methods with Random Projections}, JMLR
\textbf{21}(20):1--44 (2020), explicitly cites the source as ``Lin and
Rosasco (2017a).''  Section~4, PDF pp.~12--14, gives the answer:

\begin{itemize}
\item Theorem~2 proves optimal projected-SGM rates with no upper restriction
on its smoothness exponent $s$.
\item Corollary~5 specifies the alternative approximate-leverage-score (ALS)
Nystr\"om subsampling scheme.
\item Corollary~7 applies Theorem~2 when the projection and subsample size are
those of Corollary~5.
\item Remark~5(2) identifies the exact predecessor and says that the plain
Nystr\"om-SGM result of Lin--Rosasco (2017a) covered only
$s\in[1/2,1]$.
\end{itemize}

The notation aligns exactly.  The 2017 condition
$F_H=T^r h$ implies, for the prediction function $f_H=S_\rho F_H$,
\[
 f_H=L^{r+1/2}g,
\]
so the 2020 exponent is $s=r+1/2$.  Setting $a=0$ in Theorem~2 gives
\[
 n^{-\frac{2s}{2s+\gamma}}
 =n^{-\frac{2r+1}{2r+\gamma+1}},
\]
the same optimal rate.  The formerly saturated range $r>1/2$ is precisely
$s>1$, which Theorem~2 and Corollary~7 include.  Thus the later authors knew
the predecessor and supplied the requested different Nystr\"om subsampling
route with optimal rates.

\section*{Scope and search evidence}

The answer is in the source problem's least-squares Hilbert/RKHS setting and
under the assumptions stated in the 2020 Theorem~2 and ALS Corollary~5.  It
does not resolve the source paper's separate broad future directions on joint
cross-validation of every algorithmic parameter, unbounded kernels,
preconditioning, random features, or acceleration.

The bounded search checked the run indexes, the exact arXiv id and title,
``NySGM saturation zeta subsampling,'' ``Nystr\"om stochastic gradient
optimal rates,'' and the citation trail from arXiv:1710.07797.  The JMLR paper
is decisive because it cites the source, describes its limited smoothness
range, and extends that exact algorithmic line.

\section*{Local evidence}

\begin{itemize}
\item \texttt{source\_paper.pdf}: arXiv:1710.07797v1.
\item \texttt{supporting\_paper\_jmlr\_2020\_19-083.pdf}: JMLR 21(20),
2020.
\item Ledger: \texttt{ledger/results/1710.07797\_nysgm\_saturation\_answered\_2020.json}.
\end{itemize}

\begin{thebibliography}{9}
\bibitem{LR2017}
J. Lin and L. Rosasco,
\emph{Optimal Rates for Learning with Nystr\"om Stochastic Gradient Methods},
arXiv:1710.07797v1, 2017.

\bibitem{LC2020}
J. Lin and V. Cevher,
\emph{Convergences of Regularized Algorithms and Stochastic Gradient Methods
with Random Projections},
Journal of Machine Learning Research 21(20) (2020), 1--44.
\end{thebibliography}

\end{document}
